\documentclass[11pt,a4paper,twoside,pdf]{article}

% Encoding and typography
\usepackage[T1]{fontenc}
\usepackage[utf8]{inputenc}
\usepackage{mathpazo}

% Mathematics
\usepackage{amsmath, amssymb, amsfonts}

% Figures
\usepackage{graphicx}
\usepackage{float}
\usepackage{subcaption}

% Tables and utilities
\usepackage{multirow}
\usepackage{rotating}
\usepackage{array}
\usepackage{booktabs}
\usepackage{xspace}

% Additional symbols
\usepackage{eurosym}
\usepackage{bbding}
\usepackage{pifont}

% Other packages
\usepackage{soul}
\usepackage{color}
\usepackage{latexsym}
\usepackage{xtab}

% Hyperlinks
\usepackage[colorlinks=true,urlcolor=blue,linkcolor=blue,citecolor=blue]{hyperref}

% Equation numbering
\numberwithin{equation}{section}

% Margins
\usepackage[top=2.88cm,bottom=2.97cm,left=2.95cm,right=2.95cm]{geometry}

% Language
\usepackage[english]{babel}

% Headers
\usepackage{fancyhdr}

% Theorems
\newtheorem{theorem}{Theorem}[section]
\newtheorem{corollary}{Corollary}[theorem]
\newtheorem{lemma}[theorem]{Lemma}

% Custom commands
\newcommand{\dis}{\displaystyle}

\linespread{1.05}

\begin{document}

% Cover page %%%%%%%%%%%%%%%%%%%%%%%%%%%%%%%%%%%%%%%%%%%%%%%%%%%%%%%%%%%%%%%%%%%%%%

\pagestyle{empty}


\noindent
\begin{tabular}{r}
\includegraphics[width=8.8cm]{escudoUGRmonocromo.png} \\[-1.8ex]
\hspace{31mm}\vspace{-8mm}
\begin{tabular}{c}
\hline\\[-1ex]\hskip-2mm
{\bf Faculty of Sciences}\hspace{18mm}
\end{tabular}
\end{tabular}

\large
\vspace{30mm}
\hspace{25mm}
\begin{tabular}{l}
MASTER IN PHYSICS AND MATHEMATICS
\end{tabular}

\vspace{45mm}
\hspace{25mm}
\begin{tabular}{l}
PDEs OF TRANSPORT IN KINETIC THEORY \\
AND FLUID MECHANICS
\\[1.5ex]
\LARGE\bf VLASOV-MAXWELL
\end{tabular}
%
\vfill
\hspace{25mm}
\begin{tabular}{l}
Presented by:
\\
{\bf A. S. Amari }\\

\textcolor{blue}{\texttt{alisalemstd@gmail.com}, \texttt{alisalemstd@correo.ugr.es}}

\\[3ex]
Academic Year 2025/2026

\end{tabular}
%

\newpage
%


% Table of Contents %%%%%%%%%%%%%%%%%%%%%%%%%%%%%%%%%%%%%%%%%%%%%%%%%%%%%%%%%%%%%%%%%%%%%%%
%\newpage

\tableofcontents

% Main text %%%%%%%%%%%%%%%%%%%%%%%%%%%%%%%%%%%%%%%%%%%%%%%%%%%%%%%%%%%%%%%%%%%%%%%%
%\newpage

\pagestyle{fancy}
\fancyhead[RO,LE]{\leftmark}
\fancyhead[LO,RE]{\thepage}
\fancyfoot{}




\newpage

\section{The Vlasov--Maxwell System}

Plasmas constitute a state of matter characterized by the presence of charged particles moving at high velocities. Unlike neutral fluids, the dynamics of a plasma are governed by the joint evolution of particles and the fields they themselves generate, giving rise to a strongly coupled system. Due to the large number of particles in the system (on the order of Avogadro's number $N_A \sim 10^{23}$), the appropriate description is of a kinetic nature and is based on a probability distribution function $f(t,x,p)$ describing particle density in phase space. Additionally, an additional formulation is needed for electromagnetic interactions between particles.




%----------------------------------------------------




\subsection{Microscopic dynamics and Vlasov equation}

The microscopic description of a plasma starts from the fundamental fact that it consists of an extremely large number of charged particles (electrons and ions) that interact with each other through long-range forces. From a strictly Newtonian point of view, the state of each particle $i$ is determined by its position $x_i(t)$ and its velocity $v_i(t)$, whose evolution is given by:
\begin{equation}
\begin{cases}
x'_i(t) = v_i(t), \\[0.3em]
m\, v'_i(t) = F\big(t,x_i(t),v_i(t)\big),
\end{cases}
\label{eq:newton_micro}
\end{equation}
where $F(t,x,v)$ represents the total force field acting on the particle. This field can depend explicitly on time, position, and velocity, allowing the inclusion of both external forces and electromagnetic interactions generated by the plasma itself.

However, when the number of particles $N$ is on the order of Avogadro's number, the system \eqref{eq:newton_micro} becomes completely intractable. Even if it were possible to solve Newton's equations for each particle, the information obtained would be useless from a physical point of view: what matters is not the individual trajectory of each electron, but the collective behavior of the ensemble. This observation motivates the transition to a statistical description.

\medskip

Rather than tracking each particle individually, we introduce the distribution function


\[
f(t,x,v),
\]


which describes particle density in phase space. For any region $A\subset\mathbb{R}^3_x$ and $B\subset\mathbb{R}^3_v$, the quantity


\[
\int_A\int_B f(t,x,v)\,dx\,dv
\]


represents the number of particles that at time $t$ are in $A$ with velocities in $B$. This function is the central object of kinetic theory \cite{romeroedp}.

\medskip


To derive the equation governing the evolution of $f$, we first consider the exact discrete distribution

\begin{equation}
f_{ex}(t,x,v)=\sum_{i=1}^N \delta(x-x_i(t))\,\delta(v-v_i(t)),
\label{eq:f_ex}
\end{equation}

(where $\delta(x)$ is the Dirac delta distribution) representing the microscopic state of the system. For any smooth test function $\phi(x,v)$ with compact support, we have


\[
\int_{\mathbb{R}^3}\int_{\mathbb{R}^3} \phi(x,v)\,f_{ex}(t,x,v)\,dx\,dv
= \sum_{i=1}^N \phi(x_i(t),v_i(t)).
\]



Taking the time derivative and applying the chain rule together with \eqref{eq:newton_micro}, we obtain


\[
\frac{d}{dt}\int \phi\,f_{ex}
= \sum_{i=1}^N \left[
v_i(t)\cdot\nabla_x\phi(x_i(t),v_i(t))
+ \frac{1}{m}F(t,x_i(t),v_i(t))\cdot\nabla_v\phi(x_i(t),v_i(t))
\right].
\]



Using the Dirac delta property again, this expression can be rewritten as


\[
\int \phi\,\partial_t f_{ex}
= \int \left[
v\cdot\nabla_x\phi(x,v)
+ \frac{1}{m}F(t,x,v)\cdot\nabla_v\phi(x,v)
\right] f_{ex}(t,x,v)\,dx\,dv.
\]



\medskip

To obtain the differential equation satisfied by $f_{ex}$, we perform an integration by parts on the right-hand side. Since $\phi$ has compact support, the boundary terms vanish, and we obtain


\[
\int \left[
\partial_t f_{ex}
+ \operatorname{div}_x(v f_{ex})
+ \operatorname{div}_v\!\left(\frac{F(t,x,v)}{m} f_{ex}\right)
\right]\phi\,dx\,dv = 0.
\]



Since this equality must hold for all test functions $\phi$, $f_{ex}$ satisfies the equation in weak sense


\[
\partial_t f_{ex}
+ \operatorname{div}_x(v f_{ex})
+ \operatorname{div}_v\!\left(\frac{F(t,x,v)}{m} f_{ex}\right)=0.
\]



Finally, taking the continuous limit $f_{ex}\to f$ as $N\to\infty$, we obtain the Vlasov equation in its general conservative form:
\begin{equation}
\partial_t f + v\cdot\nabla_x f
+ \operatorname{div}_v\!\left(\frac{F(t,x,v)}{m}f\right)=0.
\label{eq:vlasov_general}
\end{equation}

\medskip


Equation \eqref{eq:vlasov_general} expresses that particle density is transported along trajectories determined by the vector field $(v,F/m)$ in phase space. It contains no diffusive or dissipative terms, so the model describes \emph{collisionless} plasmas \footnote{In this regard, a collision term can be introduced on the right-hand side of the time equation $\left( \frac{\partial f}{\partial t} \right)_c$ that expresses the rate of change of $f$ due to collisions \cite{chen1984introduction}.}. The dynamics are purely deterministic and conserve density along the flow, making the Vlasov equation a conservation law in phase space.



%------------------------------------------------------

\subsection{Electromagnetic coupling: Maxwell equations}

The Vlasov equation describes the evolution of the distribution function under the action of a general force field $F(t,x,v)$. In the case of plasmas, this field is not external but is generated by the particles of the system themselves. This introduces a fundamental coupling: particles move under the action of electromagnetic fields, and these fields in turn depend on the distribution of particles.

For charged particles, the force field is given by the Lorentz force, which in terms of the electric field $E(t,x)$ and the magnetic field $B(t,x)$ is written as \cite{martin2021campo}


\[
F(t,x,v) = q\big(E(t,x) + v \times B(t,x)\big),
\]


where $q$ is the particle charge.



The fields $E$ and $B$ are not arbitrary: their evolution is governed by Maxwell's equations, which in differential form are
\begin{align}
\nabla \cdot E &= \frac{\rho}{\varepsilon_0}, \label{eq:gauss_e}\\
\nabla \cdot B &= 0, \label{eq:gauss_b}\\
\nabla \times E &= -\frac{\partial B}{\partial t}, \label{eq:faraday}\\
\nabla \times B &= \mu_0 j + \mu_0 \varepsilon_0 \frac{\partial E}{\partial t}. \label{eq:ampere}
\end{align}
Equations \eqref{eq:gauss_e} and \eqref{eq:gauss_b} are elliptic-type equations that impose instantaneous constraints on the fields, while \eqref{eq:faraday} and \eqref{eq:ampere} are evolutionary equations that describe the propagation of electromagnetic waves.


In a plasma described by a distribution function $f(t,x,v)$, the sources of Maxwell's equations—the charge density $\rho$ and the current density $j$—are obtained as moments of $f$
\begin{equation}
\rho(t,x) = q \int_{\mathbb{R}^3} f(t,x,v)\,dv,
\qquad
j(t,x) = q \int_{\mathbb{R}^3} v\, f(t,x,v)\,dv.
\label{eq:rho_J}
\end{equation}
\\


\newpage
The coupling between the Vlasov equation and Maxwell's equations gives rise to the Vlasov--Maxwell system: the distribution function evolves under the action of electromagnetic fields, while these fields are determined by the charges and currents induced by the distribution itself. From a physical point of view, the Vlasov--Maxwell system combines two levels of description: the microscopic dynamics of particles, governed by the Vlasov equation, and the macroscopic dynamics of fields, governed by Maxwell's equations.\newline

The interaction between these two levels is what allows capturing the complexity of collisionless plasmas, where collective effects dominate over particle-particle collisions. This framework is the starting point for the derivation of reduced models, such as the one-dimensional system studied in subsequent sections.


%----------------------------------------------------------------


\section{Relativistic dynamics}




In laser-plasma interaction applications, particles can reach velocities comparable to the speed of light. Under these conditions, the classical relationship between velocity and momentum ceases to be valid, and it is necessary to resort to the relativistic formulation of dynamics \cite{janssen2022gravitacion}. The use of linear momentum $p$ as the fundamental variable is especially convenient, as it allows expressing velocity through the relativistic relation


\[
v(p)=\frac{p}{m\gamma(p)}, \qquad
\gamma(p)=\sqrt{1+\frac{|p|^2}{m^2c^2}}.
\]




Furthermore, the electromagnetic description is expressed naturally in terms of electromagnetic potentials: a scalar potential $\Phi$ and a vector potential $A$, from which the electric and magnetic fields are obtained. \newline

The objective of this section is to present the fundamental elements of relativistic dynamics of plasmas that motivate the reduced Vlasov--Maxwell model studied in this work. We introduce the relationship between momentum and velocity in the relativistic framework, then describe the role of electromagnetic potentials, and finally justify the one-dimensional reduction that leads to the non-relativistic (NR), quasi-relativistic (QR), and fully relativistic (FR) models.


%----------------------------------------------------------------


\subsection{Electromagnetic potentials}

The classical formulation of Maxwell's equations uses the electric field $E(t,x)$ and the magnetic field $B(t,x)$ as fundamental variables. However, in relativistic treatments it is more convenient to work with \emph{electromagnetic potentials}: a scalar potential $\Phi(t,x)$ and a vector potential $A(t,x)$. These potentials allow expressing the fields through the relations \cite{martin2021campo}
\begin{equation}
B = \nabla \times A, \qquad
E = -\nabla \Phi - \partial_t A.
\label{eq:potenciales}
\end{equation}


The use of electromagnetic potentials presents several advantages in deriving reduced models of Vlasov--Maxwell. It allows naturally incorporating the geometric structure of the magnetic field, especially in configurations where the field is transverse to the propagation direction. Furthermore, it facilitates dimensional reduction, since in one spatial dimension the magnetic field can be expressed through a single component of the vector potential, simplifying the force term in the Vlasov equation by expressing $E$ and $B$ directly in terms of $A$ and $\Phi$ \cite{carrillo2006global}.




%--------------------------------------------------------------



\subsection{One-dimensional reduction and relativistic models}

The Vlasov--Maxwell system in three dimensions describes with accuracy the dynamics of a collisionless plasma. However, its mathematical and computational complexity makes it necessary in many applications to derive effective models that capture the essential phenomena at reduced cost. This is the case of laser-plasma interaction, where the problem's geometry and time scales allow for simplifications \cite{ghizzo1990vlasov}.\\



In laser wave propagation events in plasmas, the dominant dynamics occur along the direction of beam propagation, denoted by $x$. Transverse phenomena—such as lateral diffusion—are much slower and of smaller amplitude. This separation of scales allows us to assume that all relevant quantities depend only on the longitudinal coordinate:


\[
f(t,x,p) \equiv f(t,x,p_x,p_\perp), \qquad
E(t,x), \qquad B(t,x).
\]


Under this hypothesis, the three-dimensional system reduces to a model in one spatial dimension, maintaining complete dependence on the momentum.



In one spatial dimension, the magnetic field can only have transverse components. Therefore, the vector potential $A$ reduces to a single component $A(t,x)$, which satisfies a wave equation coupled with the density. The longitudinal electric field is obtained from the Poisson equation, while the transverse component of the electric field is expressed as $-\partial_t A$ \cite{martin2021campo}.\\


The relationship between velocity and momentum in the relativistic framework can be simplified depending on the physical regime considered. This leads to three versions of the model \cite{janssen2022gravitacion}:

\begin{itemize}
\item \textbf{Non-relativistic model (NR):}
We assume that velocities are small compared to $c$, so


\[
\gamma = 1.
\]


The resulting system is a class of exact solutions of the non-relativistic Vlasov--Maxwell model.

\item \textbf{Quasi-relativistic model (QR):}
We maintain the relativistic relation in the transport term,


\[
\gamma_1(p)=\sqrt{1+p^2},
\]


but take $\gamma_2=1$ in the nonlinear terms coupled to the vector potential. This model captures the essential relativistic effects without introducing the full complexity of the fully relativistic case.

\item \textbf{Fully relativistic model (FR):}
We conserve the complete expression of the Lorentz factor, which in the reduced model depends on both the momentum and the vector potential:


\[
\gamma(p,A)=\sqrt{1+p^2 + A^2}.
\]


This model describes relativistic dynamics with precision, but presents much stronger nonlinear coupling.
\end{itemize}


The three models share the same kinetic structure but differ in the degree of relativity incorporated. The NR model is mathematically simpler but insufficient in high-energy regimes. The FR model is physically more accurate, but its theoretical analysis and numerical solution are considerably more complex. The QR model constitutes an intermediate compromise: it conserves the relativistic structure in the transport, maintains manageable coupling, and is widely used in laser-plasma interaction simulations.




%-------------------------------------------------------------




\section{Existence and uniqueness results}

Under the physical hypotheses stated—one-dimensionality, transverse monokineticity, and use of the vector potential—and after a change of physical units, the complete system can be written in the variables $(t,x,p)$ as \cite{carrillo2006global}:

\begin{align}
\partial_t f
+ \frac{p}{\gamma_1}\,\partial_x f
- \left( E(t,x)
+ \frac{A(t,x)}{\gamma_2}\,\partial_x A(t,x) \right)\partial_p f &= 0, \label{eq:RM1}\\[0.4em]
\partial_{tt} A - \partial_{xx} A &= - \rho_{\gamma_2}(t,x)\,A(t,x), \label{eq:RM2}\\[0.4em]
\partial_t E &= j(t,x), \label{eq:RM3}\\[0.4em]
\partial_x E &= \rho_{\mathrm{ext}}(x) - \rho(t,x). \label{eq:RM4}
\end{align}

The macroscopic quantities $\rho$, $\rho_{\gamma_2}$, and $j$ are obtained as moments of the distribution function:


\[
\{\rho,\,\rho_{\gamma_2},\,j\}(t,x)
= \int_{\mathbb{R}} \left\{ 1,\;\frac{1}{\gamma_2},\;\frac{p}{\gamma_1} \right\} f(t,x,p)\,dp.
\]




The term $\rho_{\mathrm{ext}}(x)$ represents the background ion density, which is considered stationary on the model's time scale. The system is completed with the initial conditions
\begin{equation}
f(0,x,p)=f_0(x,p), \quad
E(0,x)=E_0(x), \quad
A(0,x)=A_0(x), \quad
\partial_t A(0,x)=A_1(x).
\label{eq:IC}
\end{equation}

\subsection{Solutions via characteristics}

The Vlasov equation \eqref{eq:RM1} is a transport equation whose flow is determined by the characteristics \cite{romeroedp}


\[
(X(s),P(s)) = (X,P)(s;t,x,p),
\]


defined by the system
\begin{equation}
\frac{dX}{ds} = \frac{P(s)}{\gamma_1(P,A)}, \qquad
\frac{dP}{ds} = -E(s,X(s)) - \frac{A(s,X(s))}{\gamma_2(P,A)}\,\partial_x A(s,X(s)),
\label{eq:char}
\end{equation}
with final conditions


\[
X(t;t,x,p)=x, \qquad P(t;t,x,p)=p.
\]



A \emph{mild} solution of the reduced system is one that satisfies the Vlasov equation along the characteristics \eqref{eq:char} and that solves Maxwell's equations \eqref{eq:RM2}--\eqref{eq:RM4} in the weak sense. With this definition, we present below the results obtained in \cite{bostan2007mild} on existence and uniqueness of solutions of this type.



\begin{theorem}[Existence and uniqueness in the QR and FR cases]
Suppose that $f_0\ge 0$ and that $(1+|p|^k)f_0\in L^1(\mathbb{R}^2)$, where $k=2$ in the QR case and $k=1$ in the FR case. Assume further that there exists a function $n_0:\mathbb{R}\to\mathbb{R}_+$, non-decreasing on $\mathbb{R}_-$ and non-increasing on $\mathbb{R}_+$, such that


\[
f_0(x,p)\le n_0(p), \qquad
\int_{\mathbb{R}} |p|^k n_0(p)\,dp <\infty, \qquad
\|n_0\|_{L^\infty}<\infty.
\]



Let $\rho_{\mathrm{ext}}\in L^1(\mathbb{R})\cap L^\infty(\mathbb{R})$ be non-negative, and suppose that the initial conditions \eqref{eq:IC} satisfy


\[
E_0' = \rho_{\mathrm{ext}} - \int_{\mathbb{R}} f_0\,dp, \qquad
A_0\in W^{2,\infty}(\mathbb{R}), \qquad
A_1\in W^{1,\infty}(\mathbb{R}).
\]



Then, the system \eqref{eq:RM1}--\eqref{eq:RM4} admits a unique global solution $(f,E,A)$ such that, for all $T>0$,


\[
f\ge 0, \qquad
(1+|p|^k)f \in L^\infty([0,T];L^1(\mathbb{R}^2)),
\]




\[
E\in W^{1,\infty}([0,T]\times\mathbb{R}), \qquad
A\in W^{2,\infty}([0,T]\times\mathbb{R}).
\]


Moreover, if $f_0\in W^{1,\infty}(\mathbb{R}^2)$, then $f\in W^{1,\infty}([0,T]\times\mathbb{R}^2)$.
\end{theorem}

This result guarantees the existence of global solutions via characteristics in the QR and FR cases. In contrast, in the non-relativistic (NR) case only local-in-time existence is available, due to the lack of control over momentum growth.\\

For numerical analysis—and in particular for the study of semi-Lagrangian schemes—it is necessary to have greater regularity on the solutions. Using arguments similar to those employed in the previous theorem, we obtain:

\begin{theorem}[Enhanced regularity]
Under additional integrability and boundedness hypotheses on $f_0$ and its derivatives, and assuming greater regularity on $\rho_{\mathrm{ext}}$, $A_0$, and $A_1$, the solution $(f,E,A)$ of the reduced system satisfies


\[
(1+|p|^k)(f + |\nabla_{x,p} f|) \in L^\infty([0,T];L^1(\mathbb{R}^2)),
\]




\[
E\in W^{2,\infty}([0,T]\times\mathbb{R}), \qquad
A\in W^{3,\infty}([0,T]\times\mathbb{R}),
\]


and, if $f_0\in W^{2,\infty}$, then $f\in W^{2,\infty}$ for all $T>0$.
\end{theorem}

As a consequence, if $f_0$ has compact support, that support remains uniformly bounded for all finite time, which is essential for the analysis of numerical schemes on unbounded domains.

\medskip

These results provide the theoretical framework necessary for the study of numerical methods applied to the reduced Vlasov--Maxwell model. In particular, they guarantee the existence of well-defined characteristics and sufficiently regular fields to justify the convergence of semi-Lagrangian schemes, which will be analyzed in the following sections.



%----------------------------------------------------------------


\section{Numerical method: semi-Lagrangian scheme}

The objective of this section is to present the fundamental ideas of the semi-Lagrangian method applied to the reduced Vlasov--Maxwell model.



The Vlasov equation describes transport in phase space in which the distribution function $f$ is conserved along the characteristics. The semi-Lagrangian method exploits precisely this property: to advance the solution in time, we find the characteristic point that arrives at $(x,p)$ at time $t^{n+1}$ and evaluate the solution there at the previous time $t^n$.

This approach avoids Courant-Friedrichs-Lewy type restrictions (as appear in finite difference methods \cite{dymer2025stable}) and allows handling large displacements in phase space, making it a tool especially suited for plasmas where the range of relativistic velocities requires very small time steps to capture high gradients in momentum.



%---------------------------------------------------------------



\subsection{Characteristic approximation}

To advance the solution of the Vlasov equation in time, the semi-Lagrangian method is based on the conservation of the distribution function along the characteristics


\[
\frac{dX}{dt} = v(P), \qquad \frac{dP}{dt} = -F(t,X),
\]


where, in the quasi-relativistic case (QR),


\[
v(p) = \frac{p}{\gamma}, \qquad \gamma = \sqrt{1+p^2}, \qquad
F(t,x) = -(E(t,x) + A(t,x)\,\partial_x A(t,x)).
\]



The objective is to approximate the characteristic point backward in time,


\[
(X,P)(t-\Delta t;t,x,p),
\]


to evaluate the solution there at the previous time. Since solving this system exactly is expensive, we use a \emph{splitting} method that separates advection in space and in momentum.


The backward-in-time displacement is approximated by three steps:

\begin{enumerate}
\item \textbf{Half advection in space:}


\[
(x,p) \longmapsto \left(x - \frac{\Delta t}{2}v(p),\; p\right).
\]



\item \textbf{Full advection in momentum:}


\[
(x,p) \longmapsto \left(x,\; p - \Delta t\,F\!\left(t - \frac{\Delta t}{2},\, x\right)\right).
\]



\item \textbf{Half advection in space:}


\[
(x,p) \longmapsto \left(x - \frac{\Delta t}{2}v(p),\; p\right).
\]


\end{enumerate}

Combining these three steps, the approximation of the characteristic point backward in time is written explicitly as:


\[
X(t-\Delta t;t,x,p) \approx x
- \frac{\Delta t}{2}v(p)
- \frac{\Delta t}{2}v\!\left(p - \Delta t\,F\!\left(t-\frac{\Delta t}{2},\, x - \frac{\Delta t}{2}v(p)\right)\right)
\]
\[
P(t-\Delta t;t,x,p) \approx p
- \Delta t\,F\!\left(t-\frac{\Delta t}{2},\, x - \frac{\Delta t}{2}v(p)\right).
\]



To abbreviate, we introduce the notation


\[
\tilde{F} = F\!\left(t-\frac{\Delta t}{2},\, x - \frac{\Delta t}{2}v(p)\right),
\]

so that the previous formulas can be written as


\[
\begin{aligned}
\tilde{p} &= p - \Delta t\,\tilde{F}, \\[0.3em]
\tilde{x} &= x - \frac{\Delta t}{2}v(p) - \frac{\Delta t}{2}v(\tilde{p}).
\end{aligned}
\]



These expressions constitute the basis of the semi-Lagrangian method: once the point $(\tilde{x},\tilde{p})$ is calculated, the solution is updated via:


\[
f^{n+1}(x,p) = f^{n}(\tilde{x},\tilde{p})
\]


using interpolation to evaluate $f^n$ at non-nodal points.\\


Next we present the results obtained in \cite{bostan2009convergence} on the error in characteristic approximation and linear interpolation.

\begin{lemma}[Error in characteristic approximation]
Let $(X,P)$ be the exact characteristic and $(\tilde{x},\tilde{p})$ the approximation obtained by splitting. If the force field $F$ is sufficiently regular, then


\[
|\tilde{x}-X(t-\Delta t)| \le C\Delta t^3 + C\Delta t^2\|F-F_h\|_{L^\infty},
\]




\[
|\tilde{p}-P(t-\Delta t)| \le C\Delta t^3 + \Delta t\|F-F_h\|_{L^\infty}.
\]


\end{lemma}



%-------------------------------------------------------



The method requires evaluating $f$ at points that do not coincide with mesh nodes, so projection and interpolation operators are introduced. We use bilinear interpolation, which guarantees a second-order error in phase space:

\begin{lemma}[Linear interpolation error]

Let the mesh $\{(x_i,\, p_j) : (i,j)\in \mathbb{Z}^2\},$ where $x_i = i\,\Delta x, p_j = j\,\Delta p.$ Denote by $I : \ell^\infty(\mathbb{Z}^2) \longrightarrow L^\infty(\mathbb{R}^2),
\quad
\Pi : L^\infty(\mathbb{R}^2) \longrightarrow \ell^\infty(\mathbb{Z}^2)$ the linear interpolation and projection operators, respectively. Let $u\in W^{2,\infty}(\mathbb{R}^2)$. Then


\[
\|(I\circ\Pi)u - u\|_{L^\infty}
\le \frac{1}{8}(\Delta x^2 + \Delta p^2)
\max\left\{
\left\|\frac{\partial^2 u}{\partial x^2}\right\|_{L^\infty},
\left\|\frac{\partial^2 u}{\partial p^2}\right\|_{L^\infty}
\right\}.
\]



\end{lemma}




%-----------------------------------------------------------

\subsection{Update of electromagnetic fields}

In addition to the Vlasov equation, the reduced Vlasov--Maxwell model includes the evolution of the vector potential $A(t,x)$ and the longitudinal electric field $E(t,x)$. Both quantities are coupled to the distribution function through macroscopic moments, so their update must be done consistently with the evolution of $f$. \\

To do this we introduce the auxiliary variables


\[
U(t,x)=\partial_t A(t,x), \qquad V(t,x)=\partial_x A(t,x),
\]


which allow expressing the wave equation for $A$ as a first-order system.

The macroscopic system is then given by


\[
\partial_t E = \int_{\mathbb{R}} v(p)\,f(t,x,p)\,dp, \quad \partial_t A = U,
\]




\[
\partial_t U - \partial_x V = -A\,\rho(t,x), \quad
\partial_t V - \partial_x U = 0,
\]


where


\[
\rho(t,x)=\int_{\mathbb{R}} f(t,x,p)\,dp.
\]



The last two equations can be combined to obtain transport equations for the combinations




\[
\partial_t(U+V) - \partial_x(U+V) = -A\rho,
\]




\[
\partial_t(U-V) + \partial_x(U-V) = -A\rho.
\]




Integrating the previous equations along the transport characteristics gives us the expressions


\[
U(t,x) = \frac{1}{2}(U+V)(0,x+t) + \frac{1}{2}(U-V)(0,x-t)
\]
\[
- \frac{1}{2}\int_0^t (A\rho)(s,x+t-s)\,ds
- \frac{1}{2}\int_0^t (A\rho)(s,x-t+s)\,ds,
\]




\[
V(t,x) = \frac{1}{2}(U+V)(0,x+t) - \frac{1}{2}(U-V)(0,x-t)
\]
\[
- \frac{1}{2}\int_0^t (A\rho)(s,x+t-s)\,ds
+ \frac{1}{2}\int_0^t (A\rho)(s,x-t+s)\,ds.
\]




In the semi-Lagrangian method, macroscopic variables are updated at intermediate times $t^{n+1/2}$ and $t^{n+3/2}$. The numerical scheme uses the following discrete formulas:

\begin{itemize}
\item \textbf{Vector potential update:}


\[
A^{n+1}_i = A^n_i + \Delta t\,U^{n+1/2}_i.
\]



\item \textbf{Electric field update:}


\[
E^{n+3/2}_i = E^{n+1/2}_i + \Delta t\sum_j v(p_j)\,f^{n+1}_{ij}\,\Delta p.
\]



\item \textbf{Update of $U$ and $V$:}


\[
(U\pm V)^{n+3/2}_i =
I\big((U\pm V)^{n+1/2}\big)(x_i\pm\Delta t)
-\Delta t\,I(A^{n+1})\,I_1(\rho^{n+1})(x_i\pm\Delta t/2),
\]


where


\[
\rho^{n+1}_i = \sum_j f^{n+1}_{ij}\,\Delta p.
\]


\end{itemize}

These expressions allow updating all electromagnetic fields explicitly from the Vlasov solution at the new time step.


The following result demonstrated in \cite{bostan2009convergence} summarizes the behavior of macroscopic variables:

\begin{lemma}[Field stability]
If $A_0\in W^{2,\infty}$, $A_1\in W^{1,\infty}$, and $E_0$ is compatible with the Poisson equation, then the discrete solutions satisfy


\[
A^n,\; U^n,\; V^n,\; E^n \in W^{1,\infty},
\]


uniformly for $n\Delta t\le T$.
\end{lemma}

This result guarantees that the scheme maintains the regularity necessary so that the semi-Lagrangian approximation of the characteristics remains valid at each time step.

\medskip



%--------------------------------------------------------------

\subsection{Complete numerical scheme}

The complete procedure at each time step can be summarized in the following steps:

\begin{itemize}
\item \textbf{Reconstruction of the distribution function.}
From the nodal values $f^n_{ij}$, the bilinear interpolation operator is used to obtain a continuous approximation $I f^n$ in phase space.


\item \textbf{Characteristic approximation via splitting.}
The backward-in-time displacement is calculated using the second-order splitting scheme: it combines half advection in space, full advection in momentum, and another half advection in space.
The approximate characteristic point $(\tilde{x},\tilde{p})$ is used to evaluate the solution at the previous time.

\item \textbf{Explicit update of electromagnetic fields.}
The vector potential $A$ and the electric field $E$ are updated using their explicit formulas, while the auxiliary variables $U=\partial_t A$ and $V=\partial_x A$ are obtained by solving transport equations for $U\pm V$.


\item \textbf{Calculation of macroscopic moments.}
Once the new distribution function $f^{n+1}$ is obtained, the moments are calculated


\[
\rho^{n+1}(x_i)=\sum_j f^{n+1}_{ij}\,\Delta p, \qquad
j^{n+1}(x_i)=\sum_j v(p_j)f^{n+1}_{ij}\,\Delta p,
\]


which act as sources in the discretized Maxwell equations.
\end{itemize}

\medskip

Together, these four components give rise to a stable numerical scheme capable of handling large displacements in phase space while maintaining second-order accuracy in time.




%---------------------------------------------------------------


\subsection{Compact support and stability}

A fundamental aspect in the practical implementation of the semi-Lagrangian method is the behavior of the support of the distribution function over time. In the continuous model, if the initial condition $f_0$ has compact support in phase space, then the characteristics remain bounded for all finite time, guaranteeing that the solution $f(t)$ conserves a uniformly compact support in $t$. This fact is essential to prevent the appearance of effects without physical relation due to the truncation of the numerical domain.

In the semi-Lagrangian scheme, the approximation of characteristics via splitting and the explicit update of electromagnetic fields allow us to show that, if the discrete initial condition $(f^0_{ij})_{i,j\in\mathbb{Z}}$ has compact support, then the support of $(f^n_{ij})_{i,j\in\mathbb{Z}}$ remains uniformly bounded for all $n\Delta t\le T$. In particular, no numerical particles are generated outside the relevant physical domain. This behavior is summarized in the following result demonstrated in \cite{bostan2009convergence}:

\begin{lemma}[Compact support conservation]
If the initial condition $f_0$ has compact support and the initial fields $(E_0,A_0,A_1)$ are sufficiently regular, then there exists a compact set $K\subset\mathbb{R}^2$ such that


\[
\operatorname{supp}(f^n)\subset K \qquad \text{for all } n\Delta t\le T.
\]


\end{lemma}

Compact support conservation has two direct consequences:

\begin{itemize}
\item \textbf{Numerical stability.}
The method does not generate values without physical correlation outside the physical domain, avoiding instabilities associated with extrapolations or artificial support propagation.

\item \textbf{Computational efficiency.}
The numerical domain can be safely truncated without the need to impose artificial boundary conditions, since the solution remains confined to a finite region of phase space.
\end{itemize}

Together, compact support conservation and uniform regularity of electromagnetic fields guarantee that the semi-Lagrangian scheme is stable and appropriate for the convergence analysis developed in the next section.



%---------------------------------------------------------------
\newpage
\section{Convergence analysis}

This section presents the main results regarding the convergence of the semi-Lagrangian method applied to the reduced Vlasov--Maxwell model. The objective is to estimate the error committed by the numerical scheme in approximating the exact solution $(f,E,A)$ of the continuous system. For this, we introduce the quantities:


\[
\Delta f^n_{ij} = f^n_{ij} - f(t^n,x_i,p_j), \qquad
\Delta E^{n+1/2}_i = E^{n+1/2}_i - E(t^{n+1/2},x_i),
\]




\[
\Delta A^n_i = A^n_i - A(t^n,x_i), \qquad
\Delta U^{n+1/2}_i = U^{n+1/2}_i - \partial_t A(t^{n+1/2},x_i),
\]




\[
\Delta V^{n+1/2}_i = V^{n+1/2}_i - \partial_x A(t^{n+1/2},x_i).
\]



Next we collect the fundamental results demonstrated in \cite{bostan2009convergence} that allow us to control these errors and demonstrate the convergence of the method.\\




\begin{lemma}[Error control in $f$]
Under the hypotheses of Theorems 3.1 and 3.2, there exists a constant $C>0$ such that


\[
\|\Delta f^{\,n+1}\|_{l^\infty}
\le \|\Delta f^{\,n}\|_{l^\infty}
+ C\Delta t\big(
\|\Delta E^{\,n+1/2}\|_{l^\infty}
+ \|\Delta A^{\,n}\|_{l^\infty}
+ \|\Delta U^{\,n+1/2}\|_{l^\infty}
+ \|\Delta V^{\,n+1/2}\|_{l^\infty}
\big)
\]




\[
\qquad\qquad
+ C\big(\Delta x^2 + \Delta p^2 + \Delta t^3 + \Delta t\,\Delta x^2\big).
\]


\end{lemma}



\begin{lemma}[Error control in $A$]
Under the hypotheses of Theorems 3.1 and 3.2, there exists a constant $C>0$ such that


\[
\|\Delta A^{\,n+1}\|_{l^\infty}
\le \|\Delta A^{\,n}\|_{l^\infty}
+ \Delta t\,\|\Delta U^{\,n+1/2}\|_{l^\infty}
+ C\Delta t^3.
\]


\end{lemma}



\begin{lemma}[Error control in $E$]
Under the hypotheses of Theorems 3.1 and 3.2, there exists a constant $C>0$ such that


\[
\|\Delta E^{\,n+3/2}\|_{l^\infty}
\le \|\Delta E^{\,n+1/2}\|_{l^\infty}
+ C\Delta t\,\|\Delta f^{\,n+1}\|_{l^\infty}
+ C\Delta t\,\Delta p^2
+ C\Delta t^3.
\]


\end{lemma}


\begin{lemma}[Error control in $U\pm V$]
Under the hypotheses of Theorems 3.1 and 3.2, there exists a constant $C>0$ such that


\[
\|\Delta(U^{\,n+3/2} \pm V^{\,n+3/2})\|_{l^\infty}
\le
\|\Delta(U^{\,n+1/2} \pm V^{\,n+1/2})\|_{l^\infty}
\]




\[
\qquad\qquad
+ C\Delta t\big(\|\Delta f^{\,n+1}\|_{l^\infty}
+ \|\Delta A^{\,n+1}\|_{l^\infty}\big)
+ C\Delta x^2
+ C\Delta t(\Delta x^2 + \Delta p^2).
\]


\end{lemma}



To obtain a global measure of the error, we introduce the quantity


\[
e_n :=
\|\Delta f^{\,n}\|_{l^\infty}
+ \|\Delta A^{\,n}\|_{l^\infty}
+ \|\Delta E^{\,n+1/2}\|_{l^\infty}
\]
\[
+ \|\Delta(U^{\,n+1/2}+V^{\,n+1/2})\|_{l^\infty}
+ \|\Delta(U^{\,n+1/2}-V^{\,n+1/2})\|_{l^\infty}.
\]



The following theorem provides the final convergence estimate for the method.

\begin{theorem}[Scheme convergence]
Under the hypotheses of Theorems 3.1 and 3.2, there exists a constant $C>0$ such that, for all $n$ with $n\Delta t\le T$,


\[
e_n \le
C\big(\Delta t^2 + \Delta x^2 + \Delta p^2 + \Delta t\,\Delta p^2\big)
+ C\,\frac{\Delta x^2 + \Delta p^2}{\Delta t}.
\]


\end{theorem}


Where the last term $\frac{\Delta x^2 + \Delta p^2}{\Delta t}$ allows performing a very fine temporal discretization $\Delta t$ on the sole condition of performing discretizations $\Delta x$ and $\Delta p$ of at least half order of magnitude. This enables very high temporal precision without paying the cost of even finer spatial discretization. This is exactly the opposite of what happens in finite difference methods where the usual condition is $\Delta t \sim \left( \Delta x \right)^2$.




\section{Numerical results}

This section presents numerical results illustrating the behavior of the method and confirming the theoretical estimates obtained.\\

The code used in this work is available publicly on GitHub: \textcolor{blue}{\texttt{https://github.com/Alisama20}}.



\subsection{Reduction to 1D1V: Vlasov--Poisson and two-stream instability}



The reduced Vlasov--Maxwell model considered in this work describes the dynamics of a plasma in one spatial dimension but includes a transverse vector potential $A(t,x)$ representing the transverse electromagnetic component of the field. This potential appears coupled to the distribution function through the force


\[
F = -\left(E + A\,\partial_x A\right),
\]


and satisfies a nonlinear wave equation. However, in one velocity dimension the magnetic field vanishes \cite{martin2021campo}.

In this way, the reduced Vlasov--Maxwell system naturally collapses to the electrostatic Vlasov--Poisson model:


\[
\partial_t f + v\,\partial_x f - E(t,x)\,\partial_v f = 0,
\qquad
\partial_x E(t,x) = \rho(t,x) - \rho_0,
\]


where $\rho(t,x)=\int_{\mathbb{R}} f(t,x,v)\,dv$ and $\rho_0$ is the mean density ensuring charge neutrality. Thus, this model describes particle interaction solely through the longitudinal electric field.\\


The initial condition corresponds to two counter-streaming electron beams slightly perturbed in space:


\[
f_0(x,v) = \frac{1}{\sqrt{2\pi}}\,v^2 e^{-v^2/2}\,\big(1+\alpha\cos(\theta x)\big),
\]


with $\alpha=10^{-2}$ and $\theta=0.5$. The function $v^2 e^{-v^2/2}$ exhibits two symmetric maxima in velocity, representing two populations of particles moving in opposite directions. The small-amplitude spatial perturbation triggers the well-known \textit{two-stream instability}.

The electric field is obtained by solving the Poisson equation in Fourier space, and the Vlasov equation is integrated using the semi-Lagrangian scheme with Strang splitting (half step in velocity, full step in space, and another half step in velocity).



In Figure~\ref{doshaces}, states of the distribution function $f(t,x,v)$ in the $(x,v)$ plane are shown at times $t=0$, $t=15$, $t=20$, and $t=30$. The observed evolution is characteristic of two-stream instability:

\begin{itemize}
    \item \textbf{At $t=0$} we see the two velocity maxima corresponding to the two initial beams, with no spatial difference.

    \item \textbf{At $t\approx 15$} the linear phase of the instability begins: the spatial perturbation grows exponentially and the electric field amplifies. The distribution begins to deform in phase space.

    \item \textbf{At $t\approx 20$} fine structures and filaments appear in the $(x,v)$ plane, resulting from beam mixing and unstable mode growth.

    \item \textbf{At $t\approx 30$} electrostatic vortices and closed structures are observed. The distribution is strongly mixed and the instability has saturated.
\end{itemize}

\newpage

\begin{figure}[H]
    \centering
    \includegraphics[width=0.9\linewidth]{doshaces.png}
    \caption{The spatial domain is periodic, $x\in[0,L_x]$ with $L_x=4\pi$, discretized with $N_x=200$ points. The velocity domain is $v\in[-v_{\max},v_{\max}]$ with $v_{\max}=6$ and $N_v=200$ points. The time step is set to $\Delta t = 0.05$.}
    \label{doshaces}
\end{figure}



%---------------------------------------------------------------




\subsection{Weibel instability in 1D2V relativistic model}
\label{sec:weibel}

The model is now extended to consider two momentum components $(p_x,p_y)$ and a transverse vector potential $A_y(t,x)$ that generates a transverse magnetic field $B_z = \partial_x A_y$. The distribution function $f(t,x,p_x,p_y)$ satisfies the relativistic Vlasov equation:

\[
\partial_t f + \frac{p_x}{\gamma}\,\partial_x f
+ F_x\,\partial_{p_x} f + F_y\,\partial_{p_y} f = 0,
\]

where $\gamma = \sqrt{1 + p_x^2 + p_y^2}$ and the force on the electron ($q=-1$, $m=1$) is

\[
F_x = -E_x - \frac{p_y}{\gamma}\,B_z, \qquad
F_y = U + \frac{p_x}{\gamma}\,B_z,
\]

with $E_x$ the longitudinal electric field obtained from Poisson, $U = \partial_t A_y$ (so that $E_y = -U$), and $B_z = \partial_x A_y$ the transverse magnetic field. The vector potential $A_y$ evolves according to the wave equation coupled to the transverse current

\[
\partial_{tt} A_y - \partial_{xx} A_y = - J_y,
\quad J_y = \int \frac{p_y}{\gamma}\, f\,dp_x\,dp_y.
\]

Unlike the previous section, here we exploit the genuine electromagnetic coupling of the 1D2V model to reproduce the \textbf{Weibel instability}: when the initial distribution exhibits temperature anisotropy ($T_y > T_x$), an infinitesimal perturbation of the transverse magnetic field grows exponentially during the linear phase, saturating through particle trapping and generating current filaments in phase space. This is one of the most relevant mechanisms in plasma astrophysics and collisionless shock fronts.

To numerically integrate the system, the same semi-Lagrangian strategy as in 1D1V was followed, with symmetric Strang splitting (palindromic) of second order:

\begin{enumerate}
    \item Half step in $p_y$ and $p_x$ with cubic interpolation (Catmull--Rom) along the Lorentz force.
    \item Full step in $x$ using periodic cubic interpolation.
    \item Update of the electromagnetic field $(A_y,U)$ using an explicit Störmer--Verlet scheme:
    \[
    U^{n+1/2} = U^{n-1/2} + \Delta t\,(\partial_{xx}A_y^n - J_y^n),\qquad
    A_y^{n+1} = A_y^n + \Delta t\, U^{n+1/2}.
    \]
    \item Final half step in $p_x$ and $p_y$ using fields centered at $t^n$ to preserve splitting symmetry.
\end{enumerate}

The initial condition is an anisotropic bi-Maxwellian uniform in space, perturbed only electromagnetically through a seed in the vector potential:

\begin{equation}
f_0(x,p_x,p_y) =
\frac{1}{2\pi\sqrt{T_x T_y}}
\exp\!\left(-\frac{p_x^2}{2T_x} - \frac{p_y^2}{2T_y}\right),
\qquad A_y(x,0) = \varepsilon\cos(k_0 x),
\label{eq:CI_weibel}
\end{equation}

with parameters $T_x=1$, $T_y=4$ (anisotropy $T_y/T_x = 4$), $k_0 = 0.5$, $\varepsilon = 10^{-2}$, and $U(x,0)=0$. The domain is periodic $x\in[0,L_x]$ with $L_x = 4\pi$, discretized with $N_x = 128$ points, and the momentum domains are $p_x\in[-6,6]$ and $p_y\in[-9,9]$ with $N_{p_x} = N_{p_y} = 64$. The time step is $\Delta t = 0.02$ and the final time is $t_{\max} = 50$. The electromagnetic seed induces an initial magnetic field $B_z(x,0) = -\varepsilon k_0 \sin(k_0 x)$ on which the instability acts.

\begin{figure}[H]
    \centering
    \includegraphics[width=\textwidth]{../weibel_campos.png}
    \caption{Temporal evolution of the transverse magnetic field $\max|B_z(t)|$ (left) and the longitudinal electric field $\max|E_x(t)|$ (right) in semilogarithmic scale. The red dashed line corresponds to the exponential fit $B_0\,e^{\gamma_{\rm sim} t}$ performed on the automatically detected linear phase.}
    \label{fig:weibel_campos}
\end{figure}

In Figure~\ref{fig:weibel_campos}, we see the characteristic signature of the instability: after an initial transient, $B_z$ grows exponentially over several orders of magnitude, with a rate $\gamma_{\rm sim}$ compatible with kinetic predictions from the Weibel dispersion relation. The longitudinal electric field $E_x$ grows secondarily, fed by nonlinear coupling through the currents generated by particle rotation in the self-generated magnetic field.

\begin{figure}[H]
    \centering
    \includegraphics[width=\textwidth]{../weibel_fpx.png}
    \caption{Perturbation of the marginal distribution $\Delta f(x,p_x) = \int f\,dp_y - \int f_0\,dp_y$ at $t=0,15,30,50$. The divergent colormap highlights the spatial modulation induced by the instability.}
    \label{fig:weibel_fpx}
\end{figure}

\begin{figure}[H]
    \centering
    \includegraphics[width=\textwidth]{../weibel_fpy.png}
    \caption{Perturbation of the marginal distribution $\Delta f(x,p_y)$ at the same times. The coupling to $B_z$ preferentially deflects particles in the $p_y$ direction, visible as spatial periodic modulation in phase with the seed mode $k_0$.}
    \label{fig:weibel_fpy}
\end{figure}

Figures~\ref{fig:weibel_fpx} and~\ref{fig:weibel_fpy} show how the distribution function, initially homogeneous in $x$, develops spatial modulation with wavelength $2\pi/k_0$ reflecting the Weibel mode structure. The perturbation is larger in the $p_y$ direction, where thermal anisotropy provides the free energy feeding the instability. This coupling $(p_x,p_y)$ mediated by $B_z$ is intrinsically electromagnetic and has no equivalent in the electrostatic 1D1V model of the previous section.

%--------------------------------------------------------------

\subsection{Solution via Monte Carlo methods: PIC (Particle-in-Cell) scheme}
\label{sec:pic}

As an alternative to the semi-Lagrangian mesh scheme employed so far, this section presents a solution using \textbf{Monte Carlo methods} for the non-relativistic Vlasov--Poisson system. The method adopted is the \textbf{Particle-in-Cell (PIC)}, which constitutes the natural stochastic translation of the Vlasov equation: the distribution function is represented by a finite set of $N$ macroparticles randomly sampled from $f_0(x,v)$,
\[
f(t,x,v)\approx \sum_{p=1}^{N} w_p\,\delta\!\big(x-x_p(t)\big)\,\delta\!\big(v-v_p(t)\big),
\]
whose trajectories follow the characteristics of the transport operator, while fields are calculated on a fixed mesh. The method is Monte Carlo in the classical sense: the statistical error decays as $\mathcal{O}(N^{-1/2})$ and is independent of the phase space dimension, in contrast to the exponential cost of mesh methods.

\subsubsection*{System and algorithm}

We integrate the non-relativistic 1D1V electrostatic Vlasov--Poisson system,
\[
\partial_t f + v\,\partial_x f - E(x)\,\partial_v f = 0,\qquad
\partial_x E = 1 - \int f\,dv,
\]
with fixed background ions of unit density. Each time step consists of four stages:

\begin{enumerate}
    \item \textbf{Deposit (CIC).} Each particle deposits its weight to the two nearest mesh nodes linearly (\emph{Cloud-In-Cell}), obtaining the density $n_e(x_i)$.
    \item \textbf{Electric field.} Poisson is solved by FFT, $\hat{E}(k) = -i\,\widehat{(1-n_e)}/k$, with $\hat{E}(0)=0$.
    \item \textbf{Interpolation.} $E$ is interpolated to particle positions using the same CIC scheme (guarantees absence of self-force).
    \item \textbf{Push (leapfrog).}
    $v_p^{n+1/2} = v_p^{n-1/2} + \Delta t\,(-E_p^{\,n})$,
    $x_p^{n+1} = x_p^{n} + \Delta t\, v_p^{n+1/2}$,
    initialized using the Boris half-step retraction.
\end{enumerate}

\subsubsection*{Two-stream instability}

We reproduce the two-stream instability again, but now solved by PIC rather than with mesh in $(x,v)$. The initial distribution is the sum of two counter-propagating Gaussians with a sinusoidal spatial perturbation,
\[
f_0(x,v) = \frac{1 + \alpha \cos(k_0 x)}{2\sqrt{2\pi}\,v_T}
\left[
e^{-(v-v_0)^2/(2v_T^2)} + e^{-(v+v_0)^2/(2v_T^2)}
\right],
\]
with $v_0 = 2$, $v_T = 0.3$, $k_0 = 0.5$, $\alpha = 0.05$, and $L_x=4\pi$. These parameters satisfy the condition $k_0 v_0 < \sqrt{2}$ necessary for the mode to be unstable in the cold two-stream limit. Position sampling respecting the modulation $1+\alpha\cos(k_0 x)$ is performed through numerical inversion (Newton) of the CDF $F(x)=x+(\alpha/k_0)\sin(k_0 x)$. We use $N = 2\cdot 10^{5}$ macroparticles, $N_x=128$ mesh cells, and $\Delta t = 0.05$ until $t_{\max}=40$.

\begin{figure}[H]
    \centering
    \includegraphics[width=\textwidth]{../twostream_campos.png}
    \caption{Diagnostics of the PIC simulation: evolution of $\max|E(t)|$ (left) and electrostatic energy $\tfrac{1}{2}\int E^2\,dx$ (right). The dashed line is the exponential fit in the linear phase, compared with the theoretical rate for cold two-stream beams $\omega^2 = k_0^2 v_0^2 + 1 - \sqrt{1+4k_0^2 v_0^2}$.}
    \label{fig:pic_campos}
\end{figure}

\begin{figure}[H]
    \centering
    \includegraphics[width=\textwidth]{../twostream_fase.png}
    \caption{Phase space $(x,v)$ of PIC macroparticles at $t=10,20,30,40$. The formation of the characteristic BGK vortices associated with nonlinear trapping of particles by the self-consistent electric field is observed.}
    \label{fig:pic_fase}
\end{figure}

Figure~\ref{fig:pic_campos} shows exponential growth of the electric field during the linear phase, with rate $\gamma_{\rm sim}$ very close to the analytical prediction $\gamma_{\rm teo} = \sqrt{-\omega^2}$. Saturation occurs when the potential amplitude is comparable to the beam kinetic energy, giving rise to the vortices in phase space observed in Figure~\ref{fig:pic_fase}. These BGK vortices are the universal nonlinear signature of kinetic electrostatic instabilities, and their qualitative capture coincides with mesh methods, validating the stochastic scheme. Unlike the semi-Lagrangian method, whose fidelity scales with phase space mesh refinement, the PIC method degrades its accuracy with statistical noise that decays only as $1/\sqrt{N}$; however, its low cost per macroparticle and absence of positivity or conservation problems in velocity space make it the reference tool in realistic plasma simulations.


%--------------------------------------------------------------


\newpage

\section{Conclusion}

The kinetic dynamics of plasmas constitutes one of the most challenging problems in mathematical and computational physics. Although the Vlasov--Maxwell equations describe with precision the evolution of a collisionless plasma, their complexity necessitates the use of numerical methods capable of capturing fine-scale structure in phase space without introducing artifacts that distort the dynamics. In this work, we have developed and analyzed a high-precision semi-Lagrangian scheme for the reduced 1D1V and 1D2V versions of the model, including their quasi-relativistic formulation.

The theoretical analysis has established a clear framework for studying the error associated with the method. From local estimates on the distribution function, the vector potential, and the electric field, we have constructed a global error measure whose evolution is controlled by terms of order $\Delta t$, $\Delta x^2$, and $\Delta p^2$. The resulting convergence theorem guarantees the stability of the scheme under reasonable regularity hypotheses.

In the numerical part, two representative scenarios were considered. The first corresponds to the electrostatic Vlasov--Poisson model, where the two-stream instability was reproduced. The simulations show initial exponential growth, transition to the nonlinear regime, and formation of vortex-type structures. This case demonstrates the method's capacity to resolve fine structures in phase space.

The second scenario addresses relativistic Langmuir oscillation in the quasi-relativistic 1D2V model. The results obtained show oscillations due to the presence of the magnetic field, demonstrating the consistency of transport in two momentum dimensions. The evolution of the slices $f(t,x,p_x,p_y=0)$ reproduces the expected dynamics of a relativistic Langmuir wave, validating the model in a more demanding regime.\\

The semi-Lagrangian scheme developed is presented as an effective tool for simulating collisionless plasmas in reduced geometries.




%----------------------------------------------------------------

\newpage
% References %%%%%%%%%%%%%%%%%%%%%%%%%%%%%%%%%%%%%%%%%%%%%%%%%%%%%%%%%%%%%%%%%
%\newpage

\addcontentsline{toc}{section}{References}


\bibliographystyle{plain}
\bibliography{bibliografia}



\end{document}
